\subsection*{AI工具使用声明}

本文在论文撰写过程中使用了AI辅助工具进行文字润色和排版优化。所有模型建立、数据分析、结果解读和核心结论均由参赛队员独立完成。

\newpage

\begin{thebibliography}{99}

\bibitem{lo1998}
Lo Y M D, Tein M S C, Lau T K, et al. Quantitative analysis of fetal DNA in maternal plasma and serum: implications for noninvasive prenatal diagnosis[J]. American Journal of Human Genetics, 1998, 62(4): 768--775.

\bibitem{ashoor2013}
Ashoor G, Syngelaki A, Poon L C Y, et al. Fetal fraction in maternal plasma cell-free DNA at 11--13 weeks' gestation: relation to maternal and fetal characteristics[J]. Ultrasound in Obstetrics \& Gynecology, 2013, 42(1): 14--18.

\bibitem{wood2017}
Wood S N. Generalized Additive Models: An Introduction with R[M]. 2nd ed. Boca Raton: CRC Press, 2017.

\bibitem{pinheiro2000}
Pinheiro J C, Bates D M. Mixed-Effects Models in S and S-PLUS[M]. New York: Springer, 2000.

\bibitem{maesschalck2000}
De Maesschalck R, Jouan-Rimbaud D, Massart D L. The Mahalanobis distance[J]. Chemometrics and Intelligent Laboratory Systems, 2000, 50(1): 1--18.

\bibitem{rubinstein2016}
Rubinstein R Y, Kroese D P. Simulation and the Monte Carlo Method[M]. 3rd ed. Hoboken: Wiley, 2016.

\bibitem{revello2016}
Revello R, Sarno L, Ispas A, et al. Screening for trisomies by cell-free DNA testing of maternal blood: consequences of a failed result[J]. Fetal Diagnosis and Therapy, 2016, 40(2): 89--95.

\bibitem{kinnings2015}
Kinnings S L, Geis J A, Almasri E, et al. Factors affecting levels of circulating cell-free fetal DNA in maternal plasma and their implications for noninvasive prenatal testing[J]. Prenatal Diagnosis, 2015, 35(8): 816--822.

\bibitem{hastie2009}
Hastie T, Tibshirani R, Friedman J. The Elements of Statistical Learning: Data Mining, Inference, and Prediction[M]. 2nd ed. New York: Springer, 2009.

\bibitem{wang2019}
Wang E, Batey A, Struble C, et al. Gestational age and maternal weight effects on fetal cell-free DNA in maternal plasma[J]. Prenatal Diagnosis, 2013, 33(7): 662--666.

\end{thebibliography}
